\section{Background}
\label{sec:background}

% =========================================================================
% Stub outline. Target length for first draft: ~2.5 pages.
% This paper assumes the framework paper's background as given for the
% MDP/Bellman/approximate-Q-learning/MCTS treatment, and devotes more
% room to the hierarchical-RL lineage that the recursive contribution
% joins.
% =========================================================================

\subsection{Notation Recap from the Framework Paper}
\label{sec:background-recap}

We use the same notation as the companion framework paper. A finite
discounted MDP is a tuple $\Mset = (\Sset, \Aset, P, R, \gamma)$ with
state space $\Sset$, action set $\Aset(s)$ at each state, transition
kernel $P(s' \given s, a)$, reward $R(s, a, s')$, and discount factor
$\gamma \in [0, 1)$. The optimal action-value function $\Qstar$
satisfies the Bellman optimality equation:
\begin{equation}
\Qstar(s, a)
= \sum_{s'} P(s' \given s, a)
  \left[ R(s, a, s') + \gamma \max_{a'} \Qstar(s', a') \right].
\label{eq:bellman}
\end{equation}
The framework paper is built primarily as an approximate $Q$-learning
system \citep{sutton2018reinforcement, mnih2015human}: a parametric
action-value function $\Qpsi$ is fit by sampled Bellman targets on
observed transitions $(s, a, r, s')$ drawn from the LLM-induced
context process. MCTS appears in that system as an optional online
policy improvement operator over $\Qpsi$, in the expert-iteration
sense \citep{anthony2017thinking, silver2017alphagozero,
silver2018alphazero}: it uses sampled rollouts to produce
chance-node estimates $\Qmcts$ and a visit-count policy
$\pimcts(a \given s) \propto N(s, a)^{1/\tau}$, both of which serve as
training targets for $\Qpsi$ and, when present, for a parametric prior
$\ptheta$. The chance-node statistics converge to the inner
expectation in~\eqref{eq:bellman} as the simulation budget grows. The
framework paper develops these foundations in detail; we refer the
reader there for fuller treatment. In what follows we adopt the
framework paper's framing: $\Qpsi$ is the central learned object, and
MCTS is the sample-efficient search procedure that produces stronger
targets for it.

\subsection{The Structured Context MDP and its Recursive Schema}
\label{sec:background-context-mdp}

The framework paper instantiates the abstract MDP as a structured
context MDP and fits $\Qpsi$ over its states. A state is a tuple of
records (problem statement, facts, open and closed subgoals,
hypotheses, verifier-call log, committed answer, list of abandoned
identifiers, turn count). Each record carries a stable identifier.
Hypotheses additionally carry an optional \texttt{parent\_hypothesis}
field and a tuple of \texttt{children} identifiers; subgoals carry
analogous fields for nested sub-subgoals; verifier calls carry a
\texttt{follow\_up\_to} field. The framework paper leaves all of these
at their defaults in phase 1 (its first instantiation of the
approximate $Q$-learning system) and uses them only as a phase-2
forward-compatibility mechanism. The present paper treats them as the
primary object of study.

The framework paper's seven cognitive verbs (\textsc{Decompose},
\textsc{ComputeSubexpr}, \textsc{FormHypothesis}, \textsc{CheckCase},
\textsc{InvokeVerifier}, \textsc{Backtrack}, \textsc{CommitAnswer})
constitute the categorical action set $\Aset$. The $Q$-head $\Qpsi$
is a multi-head value function over this fixed verb set; the
verb-instantiator (a frozen LLM) supplies the continuous target of
each verb (the identifier of the hypothesis or subgoal being acted
on). The framework paper's full availability table conditions verb
availability on state contents (for example, \textsc{InvokeVerifier}
requires the presence of an unverified hypothesis); we refer the
reader there for the complete specification. The reward signal
attached to Bellman targets is verifier-anchored: terminal verifier
acceptance produces positive reward, terminal verifier rejection
produces negative reward, intermediate verifier calls produce small
shaped signals, and structural verbs receive no direct reward. This
verifier-anchored-reward-only principle is one of the framework
paper's load-bearing design commitments and shapes the recursive
treatment that follows.

The framework paper also positions MCTS explicitly as the policy
improvement operator over the approximate $Q$-learning system: $\Qpsi$
furnishes the leaf evaluation for MCTS rollouts, $\ptheta$ (when
present) furnishes the PUCT prior, and the MCTS outputs $\Qmcts$ and
$\pimcts$ furnish training targets back to $\Qpsi$ and $\ptheta$. The
algorithmic backbone is the chance-node-aware MCTS variant surveyed by
\citet{browne2012survey}. The present paper inherits this framing
throughout: when we discuss how recursive state structure interacts
with planning, MCTS enters as the sample-efficient procedure that
turns the approximate $Q$-learning loop into expert iteration, not as
a standalone search algorithm.

\subsection{Hierarchical Reinforcement Learning}
\label{sec:background-hrl}

Hierarchical RL is the literature most closely aligned with structural
recursion in deliberation MDPs. We summarize three threads.

\paragraph{The options framework.}
Sutton, Precup, and Singh \citep{sutton1999between} introduced options
as temporally extended actions: an option is a triple
$(I, \pi_o, \beta)$ with initiation set, internal policy, and
termination probability. An MDP augmented with options becomes a
semi-Markov decision process; planning over options can be exponentially
faster than planning over primitive actions when the option set is well
chosen. The option-critic architecture \citep{bacon2017option} learns
options end-to-end from reward. Our setting differs in two respects.
First, our hierarchy is implicit (parent-pointer references in a flat
state list) rather than explicit (option scoped over time). Second, our
verbs are not options: a verb like \textsc{Decompose} executes in one
decision step and does not enter a sub-policy with its own termination
condition. The analog of an option in our setting is the recursive
subtree that a \textsc{Decompose} action induces; that subtree is
planned over by MCTS acting as the policy improvement operator over
$\Qpsi$, not by a separate sub-policy.

\paragraph{MAXQ value decomposition.}
Dietterich \citep{dietterich2000hierarchical} decomposes the value
function along a task hierarchy: $V^{\pi}(i, s) = V^{\pi}(a, s) +
C^{\pi}(i, s, a)$, where $a$ is the chosen subtask of parent task $i$
and $C^{\pi}$ is the completion function for finishing $i$ after $a$
ends. The decomposition exposes shared structure (the same subtask used
under different parents) for transfer and faster learning. Our setting
admits an analogous decomposition: the value of decomposing hypothesis
$h$ into children $h_1, \ldots, h_k$ is the sum of values of the
children plus a completion term for assembling them into a proof of $h$.
We do not impose this decomposition architecturally; the $Q$-head
$\Qpsi$ sees the whole serialized state, not the parent-child
structure. The MAXQ framing is useful when reasoning about expected
sample efficiency of the approximate $Q$-learning system: a productive
decomposition is precisely one whose MAXQ-style completion can be
computed from verifier-anchored leaf signals.

\paragraph{Feudal and goal-conditioned hierarchies.}
Feudal networks \citep{vezhnevets2017feudal} factor a policy into a
manager that proposes goals and a worker that achieves them, with the
manager operating at a coarser timescale. Goal-conditioned hierarchical
RL \citep{nachum2018data} learns goal embeddings end-to-end. The
analogous decomposition in our setting is between subgoal generation
(via \textsc{Decompose}) and verifier-anchored evaluation of leaf
hypotheses. Unlike feudal networks, the manager and worker share the
same frozen LLM in our architecture; the hierarchy is in the state, not
in the policy.

\paragraph{Common limitation.}
Across all three threads, hierarchy is treated as an additional
architectural ingredient: separate option policies, a value-function
decomposition with cross-subtask coupling, or a two-level
manager-worker split. The companion framework paper deliberately avoids
this. The hierarchy in our setting lives entirely in the state
representation (as identifier references), the action set stays
primitive (the same seven verbs), and the policy and $Q$-heads operate
on the flat serialized state without explicit awareness of the implicit
tree. The question this paper develops is whether such an implicit
hierarchy actually purchases the benefits the explicit-hierarchy
literature claims, within an approximate $Q$-learning system whose
targets come from MCTS-improved rollouts.

\subsection{Tree-Structured Encoders and Why We Avoid Them}
\label{sec:background-trees}

Tree-LSTMs \citep{tai2015improved} and the broader recursive neural
network tradition \citep{socher2013recursive} build tree-aware encoders
by recursive composition along a parsed structure. Graph neural networks
\citep{velickovic2018gat} generalize this to arbitrary graphs via
message passing. Set transformers \citep{lee2019set} provide
permutation-invariant attention. Each comes with a representational
prior: respect tree composition, respect graph topology, respect
permutation symmetry.

For our setting, the framework paper's design choice is to embed states
via a frozen LLM operating on a serialized token sequence; the $Q$-head
$\Qpsi$ reads off the LLM's terminal hidden state. A tree-aware encoder
would require unfreezing the LLM (defeating the architectural
separation), bolting on a separate tree encoder (introducing
representation drift), or producing structured features that the LLM
cannot consume. We instead let causal self-attention resolve the tree
from depth markers and parent identifiers in the serialized prefix.
Whether this is a sufficient substitute for explicit tree encoders is
an open empirical question; Section~\ref{sec:method-serialization}
develops the argument that it should be, and
Section~\ref{sec:experimental-plan}'s ablation provides the test.
